% ============================================================
% SECTION VI: EXPERIMENTAL SETUP AND RESULTS
% ============================================================

\section{Experimental Setup and Results}

\subsection{Implementation and Training Details}
The system was implemented using Python 3.13 and TensorFlow 2.19. Model training and evaluation were performed using 5-Fold Stratified Cross-Validation on the 552 CTU-UHB recordings to ensure robust out-of-sample generalization. 

Due to the severe class imbalance (only 7.25\% of cases exhibit pathological acidemia defined by $\text{pH} < 7.05$), we utilized a Time-Series Generative Adversarial Network (TimeGAN) \cite{yoon2019} with a Wasserstein-GP objective. Crucially, to prevent data leakage, TimeGAN was trained strictly per-fold, synthesizing minority class traces exclusively from the training split of each fold. The networks were optimized using Adam with Focal Loss \cite{lin2017} ($\alpha = 0.25, \gamma = 2.0$) to heavily penalize misclassifications of the critical minority class.

\subsection{Performance Metrics and Benchmarking}
To rigorously benchmark against existing literature, we evaluated the ensemble primarily on the Area Under the Receiver Operating Characteristic Curve (AUC-ROC), alongside calibration metrics (Expected Calibration Error). 

As a further robustness check, we directly reproduced the architecture of Mendis et al.\ \cite{mendis2023} (a 1D-ResNet coupled with tabular data) and ran it through our exact V6.0 data processing pipeline to establish a fair, apples-to-apples baseline.

As summarized in Table~\ref{tab:results}, the V6.0 Tri-Modal Stacking Ensemble achieves an AUC-ROC of \textbf{0.8566} under strict, leak-free 5-Fold Cross-Validation, outperforming both the 1D-CNN baseline and the reproduced Mendis architecture.

\begin{table}[htbp]
\centering
\caption{Performance Comparison on the CTU-UHB Dataset (5-Fold CV)}
\label{tab:results}
\resizebox{\columnwidth}{!}{%
\begin{tabular}{|l|c|}
\hline
\textbf{Model Structure} & \textbf{AUC-ROC} \\ \hline
1D-CNN Baseline \cite{spilka2016} & 0.5600 \\ \hline
Fusion ResNet (Mendis 2023 Reported) \cite{mendis2023} & 0.8400 \\ \hline
Mendis Architecture (Reproduced on V6.0 Data) & 0.7983 \\ \hline
\textbf{NeuroFetal V6.0 (Stacking Ensemble)} & \textbf{0.8566} \\ \hline
\end{tabular}
}
\end{table}

\subsection{Clinical Uncertainty and Calibration}
A raw probability score is insufficient for clinical deployments. In V6.0, we applied Temperature Scaling and Platt Scaling (CalibratedClassifierCV) \cite{platt1999} to the ensemble meta-learner. The calibrated system achieves a mean Expected Calibration Error (ECE) of 0.2417 across folds, confirming improved reliability of output probabilities.

Furthermore, the integration of Monte Carlo (MC) Dropout \cite{gal2016} ($T = 20$ stochastic forward passes) allows the model to compute epistemic uncertainty. The mean predictive uncertainty (ensemble standard deviation) across evaluated windows was $\sigma = 0.3035$. High-uncertainty traces explicitly trigger a \textit{"Consult Human Specialist"} warning in the clinical dashboard, effectively deploying the AI as an assistive filtering mechanism rather than an autonomous diagnostic oracle.

\subsection{Edge Deployment Optimization}
To facilitate deployment in rural, low-resource clinical settings lacking robust internet infrastructure, the 27 MB Keras ensemble was optimized using TensorFlow Lite (TFLite). Applying Post-Training Integer (Int8) Quantization, the model footprint was compressed to \textbf{1.9 MB}.

Inference testing on commodity Android smartphone hardware demonstrated that processing a 20-minute multi-modal fetal trace takes less than 30 milliseconds on a standard mobile CPU, consuming negligible battery power. This offline capability directly operationalizes the proposed tri-modal architecture for democratized global health access.

% ============================================================
% SECTION VII: DISCUSSION & CONCLUSION
% ============================================================

\section{Discussion and Limitations}
The empirical results indicate that NeuroFetal AI outperforms classical machine learning \cite{spilka2012}, generic 1D-CNNs \cite{spilka2016}, and early fusion ensemble models \cite{mendis2023} on the CTU-UHB database. 

This performance improvement is primarily driven by the tri-modal Cross-Modal Attention Fusion mechanism. Extensive ablation studies executed during development confirmed that networks evaluating FHR isolated from UC contraction dynamics failed to distinguish pathological late decelerations from benign early decelerations. By spatially combining FHR and UC via CSP, and temporally correlating them through attention gated by clinical context, the model effectively discriminates between these distinct physiological patterns.

\subsection{Limitation: Dataset Size and Generalization}
Despite rigorous 5-Fold Cross-Validation and the correction of data leakage via per-fold TimeGAN, a critical limitation remains the size of the open-access CTU-UHB dataset (552 cases, 40 minority). While TimeGAN artificial synthesis successfully bolstered the training distribution and mitigated class imbalance, synthetic data augmentation may not fully capture entirely novel clinical edge-cases.

Recent literature demonstrates that AI models trained exclusively on specific regional fetal datasets (such as CTU-UHB) often suffer performance degradation when applied to diverse cross-institutional databases \cite{lopes2025}. Thus, while NeuroFetal V6.0 is internally calibrated and performs well on the baseline domain, the proposed architecture requires extensive multi-center clinical validation to assess its generalization across diverse demographic populations.

\section{Conclusion and Future Work}
In this paper, we proposed NeuroFetal AI, an advanced intrapartum clinical decision support system utilizing a Tri-Modal Stacking Ensemble. By fusing FHR waveforms, UC dynamics, and maternal clinical demographics via a novel Cross-Modal Attention mechanism, the model achieves a state-of-the-art, leak-free AUC-ROC of \textbf{0.8566} under rigorous 5-Fold Cross-Validation on the CTU-UHB dataset, outperforming established benchmarks.

The system facilitates practical clinical testing by integrating Monte Carlo Dropout for epistemic uncertainty quantification and TFLite (Int8) quantization, compressing the entire diagnostic pipeline into a \textbf{1.9 MB} offline edge model capable of running on low-resource mobile hardware in under 30 milliseconds.

Future iterations of NeuroFetal AI will target federated learning schemas for cross-institutional dataset integration preserving patient privacy, aiming to solidify its viability as a dependable, globally accessible tool to reduce intrapartum fetal mortality.




